% ============================================================================
% Fichier : annexes/Ann_F_contacts.tex
% Titre   : Contacts et Ressources
% Version : 1.0 -- novembre 2026
% ============================================================================
% Adresses de sites institutionnels données à titre indicatif, à vérifier
% avant diffusion. Contact de NDAP SUKULU : celui fixé par l'auteur.
% ============================================================================

\chapter{Contacts et Ressources}
\label{ann:contacts}

Les adresses ci-dessous sont celles des sites institutionnels à la date de
rédaction~; elles peuvent changer et sont à vérifier avant usage.

\section*{Autorités camerounaises et régionales}

\begin{tableaucours}{L{3.6cm}L{4.6cm}L{4.0cm}}{%
    \tableauHeader{Institution} &
    \tableauHeader{Rôle dans le cours} &
    \tableauHeader{Site}
}
ANTIC & Sécurité des SI, audits de sécurité (loi n°~2010/012) &
\texttt{www.antic.cm} \\
ART & Régulation des communications électroniques & \texttt{www.art.cm} \\
MINPOSTEL & Politique du secteur~; destinataire des rapports d'audit de
sécurité & \texttt{www.minpostel.gov.cm} \\
COBAC (Secrétariat général) & Contrôle interne et risques des
établissements de crédit de la CEMAC & \texttt{www.sgcobac.org} \\
OHADA & Droit des sociétés, responsabilité des dirigeants, comité
d'audit & \texttt{www.ohada.org} \\
Autorité de protection des données & Loi n°~2024/017 & · \\
Chambre des comptes de la Cour suprême & Contrôle des finances publiques
& · \\
Union africaine & Convention de Malabo & \texttt{au.int} \\
\end{tableaucours}

\section*{Organisations professionnelles et normatives}

\begin{tableaucours}{L{3.6cm}L{4.6cm}L{4.0cm}}{%
    \tableauHeader{Organisation} &
    \tableauHeader{Ressources} &
    \tableauHeader{Site}
}
ISACA & COBIT, certifications (CISA, CGEIT) & \texttt{www.isaca.org} \\
PeopleCert (AXELOS) & ITIL et ses certifications & \texttt{www.peoplecert.org} \\
ISO & ISO/IEC~38500, 20000-1, 27001, 27002, 31000 & \texttt{www.iso.org} \\
IIA & Normes mondiales d'audit interne, modèle des trois lignes &
\texttt{www.theiia.org} \\
COSO & Contrôle interne, gestion des risques & \texttt{www.coso.org} \\
INTOSAI & ISSAI, GUID~5100 & \texttt{www.intosai.org} \\
NIST & Définition du nuage, AI RMF & \texttt{www.nist.gov} \\
UNESCO & Recommandation sur l'éthique de l'IA & \texttt{www.unesco.org} \\
\end{tableaucours}

\section*{NDAP SUKULU}

Les fiches d'auto-évaluation et les dix labs du cours sont disponibles sur
NDAP SUKULU, sous le code \codeUE. Leur présentation et le tableau des labs
figurent au début du livre, dans « NDAP SUKULU, la salle de classe compagnon ».
NDAP SUKULU est une propriété de MaletYon, mise à la disposition des
étudiants et des praticiens souhaitant se perfectionner~:
\texttt{maletyon@proton.me}.
